\documentclass[border=14pt]{standalone}
\usepackage{ctex}
\usepackage{tikz}
\usepackage{pgfplots}
\pgfplotsset{compat=1.18}
\begin{document}
\begin{tikzpicture}[x=1cm, y=1cm, font=\normalsize, line width=0.8pt]

% fig-ge-07 555 单稳态的 RC 充电曲线（中文）
% 上行：触发脉冲（窄脉冲）；下行：电容电压从 0 升到 2/3 Ucc，标 Tp 与 1.1RC
% 说明脉宽与 Tp 无关

% ── 上行：触发脉冲 ──
\begin{scope}[shift={(0,3.6)}]
  \node[anchor=east, font=\small] at (-0.25,0.45) {触发};
  \draw[->] (0,-0.3) -- (9.2,-0.3) node[anchor=west, font=\small] {t};
  % 高电平基线
  \draw (0,0.8) -- (2.2,0.8);
  % 窄脉冲（低有效）：从 0.8 掉到 0 再回来
  \draw (2.2,0.8) -- (2.2,0) -- (2.9,0) -- (2.9,0.8);
  \draw (2.9,0.8) -- (9.0,0.8);
  \node[anchor=south, font=\small] at (2.55,1.05) {窄脉冲};
\end{scope}

% ── 下行：电容电压 ──
\begin{scope}[shift={(0,0)}]
  \node[anchor=east, font=\small] at (-0.25,0.9) {Uc};
  \draw[->] (0,0) -- (9.2,0) node[anchor=west, font=\small] {t};
  \draw[->] (0,0) -- (0,2.4) node[anchor=south, font=\small] {Uc};
  % 充电曲线：从 t0=2.9 开始，趋近 Ucc=2.4，在 2/3 Ucc=1.6 处取 Tp
  % 用指数：Uc = Ucc*(1-exp(-(t-t0)/tau))，tau=0.95 → 1/3 处 t-t0 = 1.1*tau
  \draw[domain=2.9:9.0, smooth, samples=120, line width=1pt]
    plot (\x, {2.4*(1-exp(-(\x-2.9)/0.95))});
  % 2/3 Ucc 水平虚线
  \draw[dashed, gray, line width=0.5pt] (2.9,1.6) -- (9.0,1.6);
  \node[anchor=east, font=\small] at (2.85,1.6) {2/3 Ucc};
  % Tp 标注：从 t0 到 2/3 处（t = 2.9 + 1.1*0.95 = 3.945）
  \draw[<->, line width=0.6pt] (2.9,0.15) -- (3.95,0.15);
  \node[anchor=west, font=\small] at (4.05,0.15) {Tp ≈ 1.1RC};
  \draw[dashed, gray, line width=0.5pt] (3.95,0.35) -- (3.95,1.6);
  % Ucc 参考
  \draw[dashed, gray, line width=0.5pt] (0,2.4) -- (9.0,2.4);
  \node[anchor=east, font=\small] at (-0.05,2.4) {Ucc};
  % 文字区（左下空白，远离曲线）
  \node[anchor=north west, font=\small] at (0.1,-0.55)
    {脉宽与 Tp 无关（触发脉冲只负责起跑）};
\end{scope}

\end{tikzpicture}
\end{document}